\chapter*{Appendix B: Course Source and Tutorial Map}
\addcontentsline{toc}{chapter}{Appendix B: Course Source and Tutorial Map}
\markboth{Appendix B: Source Map}{}

These notes were built from the supplied tutorial documents, the Odoo West End Bicycles build guide, the supplied forecast/S\&OP spreadsheet, and the provided third edition of \emph{Concepts in Enterprise Resource Planning}. The table below shows how the materials were aligned.

\begin{erptable}
\begin{tabularx}{\linewidth}{@{}>{\bfseries\leavevmode\color{ERPNavy}}p{0.15\linewidth}p{0.31\linewidth}X@{}}
\toprule
\rowcolor{ERPPaleBlue!92}
Tutorial & Main activity & Notes / textbook support \\
\midrule
T1 & Business functions, information flows, process design & Chapter 1 of the textbook; Odoo guide provides West End Bicycles organisation and the full integrated process. \\
T2 & ERP benefit and implementation challenge & Chapter 2 plus implementation/change-management concepts from Chapter 7. \\
T3 & Production approach, forecast, S\&OP & Chapter 4 production-planning hierarchy; Odoo BOM and Manufacturing Order steps. \\
T4 & CRM proposal for Luna Electronics & Chapter 3 CRM section; Odoo CRM pipeline and quotation conversion. \\
T5 & Sales Order Process / Fitter Snacker questions & Chapter 3 Fitter Snacker sales problems, ERP sales order processing, and Sales/Accounting integration. \\
T6 & Footloose forecast and S\&OP spreadsheet & Chapter 4 forecasting and S\&OP formulas; supplied spreadsheet values. \\
T7 & SCM research and CRM-vs-SCM comparison & Chapter 4 traditional/integrated supply chain and SCM measures of success; Odoo Purchase/Inventory flow. \\
T8 & HR recruitment and process modelling & Chapter 6 recruiting/hiring problems + Chapter 7 flowcharts/swimlanes; Odoo Employees organisation records. \\
\bottomrule
\end{tabularx}
\end{erptable}

\section*{Source discrepancies preserved in the notes}

The materials contain a few inconsistencies that students should recognise rather than have silently corrected for them:

\begin{itemize}
  \item The Tutorial 6 Word file asks students to forecast 2023 from 2022 sales, while the supplied spreadsheet headings still refer to 2022.
  \item Tutorial 6 states safety stock = 100; the spreadsheet pre-fills 500 in the initial December inventory cell.
  \item Tutorial 6 describes sales in pairs of shoes but production capacity as 40 shoes per hour, creating a unit-of-measure ambiguity.
  \item Tutorial numbering in some source documents is inconsistent (for example, a page header may use a different tutorial number from the title). The notes organise topics by subject rather than reproducing those numbering mismatches.
\end{itemize}

\begin{keyidea}
When a source file and a spreadsheet disagree, do not hide the conflict. State the assumption used, keep units explicit, and confirm assessment expectations with the tutor.
\end{keyidea}
